\documentclass[../main.tex]{subfiles}

\graphicspath{{../image/}{../figure/}}

\begin{document}

\chapter{标量守恒律的 L$^{1}$ 理论}\label{chapter:ch06-标量守恒律L1理论}

若干空间维数中标量平衡律的理论已达到近乎完备的境地。在经典解的框架下，特征线这一初等而有效的方法给出了\cref{theorem:ch05-凸熵局部存在}的一个更精细的版本，它显式地确定了初值 Lipschitz 连续时解的寿命区间，从而表明该寿命一般是有限的。因此，即使初值非常光滑，也必须处理弱解。


关于弱解，使标量平衡律区别于多于一个方程的方程组的特殊之处，在于其熵族的规模。下文将表明，熵的丰沛性诱导出对可容许弱解的有效刻画，以及极强的 L$^{1}$ 稳定性与 L$^{\infty}$ 单调性性质。凭借如此有力的先验估计，可以用多种方式构造可容许弱解。作为示例，此处介绍黏性消失法、L$^{1}$ 压缩半群理论、分层法、松弛法以及由动力论启发的途径。黏性消失法还将用于求解初边值问题。当初值为局部有界变差时，解亦然。然而值得注意的是，即使仅仅属于 L$^{\infty}$ 的解，也呈现出与 $BV$ 函数相同的几何结构，其跳跃间断组装在余维一的“流形”上。

本章末尾将描述在研究若干空间维数中双曲守恒律方程组弱解时所遇到的看似不可逾越的障碍，并介绍当前为绕过这些障碍所做的努力。

为展示理论的优雅，讨论将限于齐次标量守恒律；尽管一般的非齐次平衡律\eqref{align:ch03-标量平衡律}可用同一方法处理，只是伴随些许技术上的复杂化。


弱解相对于 L$^{\infty}$ 的弱$^{*}$拓扑的稳定性问题将在第 XVI 章讨论。单空间变量的特殊情形 $m=1$ 本身拥有极为丰富的理论，其中某些方面将在后续章节，特别是第 XI 章中介绍。

\subfile{Chapter-06/section-01.tex}
\subfile{Chapter-06/section-02.tex}
\subfile{Chapter-06/section-03.tex}
\subfile{Chapter-06/section-04.tex}
\subfile{Chapter-06/section-05.tex}
\subfile{Chapter-06/section-06.tex}
\subfile{Chapter-06/section-07.tex}
\subfile{Chapter-06/section-08.tex}
\subfile{Chapter-06/section-09.tex}
\subfile{Chapter-06/section-10.tex}
\subfile{Chapter-06/section-11.tex}

\end{document}
